% !TEX TS-program = pdflatex
% !TEX encoding = UTF-8 Unicode

% This file is a template using the "beamer" package to create slides for a talk or presentation
% - Giving a talk on some subject.
% - The talk is between 15min and 45min long.
% - Style is ornate.

% MODIFIED by Jonathan Kew, 2008-07-06
% The header comments and encoding in this file were modified for inclusion with TeXworks.
% The content is otherwise unchanged from the original distributed with the beamer package.

\documentclass{beamer}
\setbeamercovered{invisible}



% Copyright 2004 by Till Tantau <tantau@users.sourceforge.net>.
%
% In principle, this file can be redistributed and/or modified under
% the terms of the GNU Public License, version 2.
%
% However, this file is supposed to be a template to be modified
% for your own needs. For this reason, if you use this file as a
% template and not specifically distribute it as part of a another
% package/program, I grant the extra permission to freely copy and
% modify this file as you see fit and even to delete this copyright
% notice. 


\mode<presentation>
{
  \usetheme{Warsaw}
  % or ...

  % or whatever (possibly just delete it)
}


\usepackage[english]{babel}
% or whatever

\usepackage[utf8]{inputenc}
% or whatever

\usepackage{times}
\usepackage[T1]{fontenc}
% Or whatever. Note that the encoding and the font should match. If T1
% does not look nice, try deleting the line with the fontenc.

%%% MATH RELATED

\input{macros.tex}

\title{MATH 350-2 Advanced Calculus} 
\subtitle
{} % (optional)

\author[W.R. Casper] % (optional, use only with lots of authors)
{W.R. Casper}
% - Use the \inst{?} command only if the authors have different
%   affiliation.

\institute[California State University Fullerton] % (optional, but mostly needed)
{
  Department of Mathematics\\
  California State University Fullerton}
% - Use the \inst command only if there are several affiliations.
% - Keep it simple, no one is interested in your street address.

\subject{Talks}
% This is only inserted into the PDF information catalog. Can be left
% out. 



% If you have a file called "university-logo-filename.xxx", where xxx
% is a graphic format that can be processed by latex or pdflatex,
% resp., then you can add a logo as follows:

% \pgfdeclareimage[height=0.5cm]{university-logo}{university-logo-filename}
% \logo{\pgfuseimage{university-logo}}



% Delete this, if you do not want the table of contents to pop up at
% the beginning of each subsection:
\AtBeginSubsection[]
{
  \begin{frame}<beamer>{Outline}
    \tableofcontents[currentsection,currentsubsection]
  \end{frame}
}


% If you wish to uncover everything in a step-wise fashion, uncomment
% the following command: 

%\beamerdefaultoverlayspecification{<+->}


\begin{document}

\begin{frame}
  \titlepage
\end{frame}

\begin{frame}{Outline}
  \tableofcontents
  % You might wish to add the option [pausesections]
\end{frame}

% Since this a solution template for a generic talk, very little can
% be said about how it should be structured. However, the talk length
% of between 15min and 45min and the theme suggest that you stick to
% the following rules:  

% - Exactly two or three sections (other than the summary).
% - At *most* three subsections per section.
% - Talk about 30s to 2min per frame. So there should be between about
%   15 and 30 frames, all told.

\section{Real Analysis Lecture 8}

\subsection{Open Balls and Open Sets}

\begin{frame}{Euclidean space}
\pause
$n$-dimensional \textbf{euclidean space} is
\pause
$$\mathbb{R}^n = \{(a_1,a_2,\dots,a_n): a_1,\dots, a_n\in\mathbb{R}\}.$$
\pause
Definitions:\\
\pause
Given $\vec x = (x_1,\dots, x_n)$ and $\vec y = (y_1,\dots, y_n)$
\begin{enumerate}[(a)]
\pause
\item $\vec x + \vec y = (x_1+y_1,\dots,x_n+y_n)$
\pause
\item $a\vec x = (ax_1,\dots, ax_n)$
\pause
\item $\vec x - \vec y = \vec x + (-1)\vec y = (x_1-y_1,\dots,x_n-y_n)$
\pause
\item $\vec 0 = (0,\dots, 0)$
\pause
\item $\vec x\cdot\vec y = \sum_{i=1}^n x_iy_i$
\pause
\item $|\vec x| = (\vec x\cdot\vec x)^{1/2} = \left(\sum_{i=1}^nx_i^2\right)^{1/2}$
\end{enumerate}
\end{frame}

\begin{frame}{Metric properties}
\pause
The norm $|\vec x|$ is an example of a \vocab{metric}.\\
\pause
It satisfies several important properties:
\pause
\begin{thm}
\begin{enumerate}[(a)]
\pause
\item (positivity) $|\vec x| \geq 0$ with equality iff $\vec x = \vec 0$
\pause
\item (symmetry) $|\vec x + \vec y| = |\vec y + \vec x|$
\pause
\item (triangle inequality) $|\vec x + \vec y| \leq |\vec x| + | \vec y|$
\end{enumerate}
\end{thm}
\pause
It also satisfies 
\pause
\begin{enumerate}[(a)]
\pause
\item (scaling)  $|c\vec x| = |c|\ |\vec x|$
\pause
\item (Cauchy-Schwartz) $|\vec x\cdot\vec y|\leq |\vec x|\ |\vec y|$
\end{enumerate}
\end{frame}

\begin{frame}{Open sets}
An \textbf{open ball} of radius $r$ centered at $\vec a$ is
\pause
$$B(\vec a;r)  = \{\vec x: |\vec x-\vec a| < r\}.$$
\pause
\begin{defn}
\pause
A point $\vec a$ in a subset $A\subseteq\mathbb{R}^n$ is called an \textbf{interior point} of $A$ if there exists $r>0$ with
$B(\vec a; r)\subseteq A$.
\pause
If every point of $A$ is an interior point, then $A$ is called an \textbf{open set}.
\end{defn}
\end{frame}

\begin{frame}{Challenge!}
\begin{prob}
Prove that the empty set $\varnothing$ and the whole space $\mathbb{R}^n$ are open.
\end{prob}
\end{frame}

\begin{frame}{Challenge!}
\begin{prob}
Let $\vec a\in\mathbb{R}^n$.
Show that the singleton set
$$A = \{\vec a\}$$
is not open.
\end{prob}
\end{frame}


\begin{frame}{Challenge!}
\begin{prob}
Let $a,b,c,d\in\mathbb{R}$ with $a < b$ and $c < d$.
Prove that the \textbf{open rectangle}
$$(a,b)\times (c,d) = \{(x,y): a < x < b,\ \ c < y < d\}$$
is an open set
\end{prob}
\end{frame}

\begin{frame}{Challenge!}
\begin{prob}
Prove that an open ball is an open set.
\end{prob}
\end{frame}

\end{document}


